\partieMat{1}{Plus, moins, autant}{Comparer des quantités sans compter.}{\pic[scale=2.6] at (16.8,6.4) {panier};\pic[scale=1.5,mat/c=matRouge] at (16,7.4) {pomme};\pic[scale=1.5,mat/c=matVert] at (17.5,7.5) {pomme};\pic[scale=1.5,mat/c=matRouge] at (16.7,8.4) {pomme};\pic[scale=1.9,mat/c=matRouge] at (14,9.6) {pomme};\pic[scale=1.4,mat/c=matVert] at (18.8,9.8) {pomme};}{Avant quatre ans, l’enfant compare globalement des collections très différentes et utilise les mots « plus que », « moins que ». Il apprend aussi à comparer par correspondance terme à terme, même pour des collections plus grandes que les nombres qu’il connaît.}{\ligneAct{1}{Beaucoup ou peu ?}{act-ps-1}\ligneAct{2}{Le moins d’oiseaux}{act-ps-2}\ligneAct{3}{Une carotte pour chaque lapin}{act-ps-3}\ligneAct{4}{Il manque un gâteau}{act-ps-4}\ligneAct{5}{Une chaise pour chaque ours}{act-ps-5}\ligneAct{6}{Éléphants et fourmis}{act-ps-6}}
\begin{activite}{1}{Beaucoup ou peu ?}{entourer}{Entoure le panier où il y a {\otBold\color{otNavy}beaucoup} de pommes.}{Comparer globalement deux quantités très différentes et dire « plus que », « moins que ».}{act-ps-1}
\matFigurePage{ps-01}
\end{activite}
\begin{activite}{2}{Le moins d’oiseaux}{entourer}{Entoure l’arbre où il y a le {\otBold\color{otNavy}moins} d’oiseaux.}{Comparer globalement deux quantités très différentes et dire « moins que ».}{act-ps-2}
\matFigurePage{ps-02}
\end{activite}
\begin{activite}{3}{Une carotte pour chaque lapin}{relier,dire}{Relie chaque lapin à une carotte. Y a-t-il une carotte pour chaque lapin ?}{Comparer deux collections par correspondance terme à terme.}{act-ps-3}
\matFigurePage{ps-03}
\end{activite}
\begin{activite}{4}{Il manque un gâteau}{relier,dessiner}{Relie chaque gâteau à une assiette. Dessine le gâteau qui manque.}{Comparer par correspondance terme à terme, même avec des collections plus grandes que les nombres connus.}{act-ps-4}
\matFigurePage{ps-04}
\end{activite}
\begin{activite}{5}{Une chaise pour chaque ours}{relier,entourer}{Relie chaque ours à une chaise. Entoure l’ours qui n’a pas de chaise.}{Comparer deux collections par correspondance terme à terme.}{act-ps-5}
\matFigurePage{ps-05}
\end{activite}
\begin{activite}{6}{Éléphants et fourmis}{relier,dire}{Relie chaque éléphant à une fourmi. Y a-t-il autant d’éléphants que de fourmis ?}{Comprendre qu’une quantité ne dépend ni de la taille ni de la place des objets.}{act-ps-6}
\matFigurePage{ps-06}
\end{activite}
\partieMat{2}{Un, deux, trois}{Découvrir les premiers nombres.}{\pic[scale=1.7,mat/n=3] at (14.6,9) {de};\pic[scale=1.7,mat/n=2] at (17.2,7) {de};\pic[scale=2.1,mat/n=3] at (18.6,9.6) {main};\pic[scale=1.3,mat/n=1] at (13.2,6.6) {de};}{Les petits nombres s’apprennent dans l’ordre : deux, puis trois, voire quatre. L’enfant les rencontre avec ses doigts, des constellations de points et des collections variées, avant de découvrir les chiffres. On ne parle pas encore de zéro.}{\ligneAct{7}{Un seul}{act-ps-7}\ligneAct{8}{Deux}{act-ps-8}\ligneAct{9}{Trois}{act-ps-9}\ligneAct{10}{Un de plus}{act-ps-10}\ligneAct{11}{Mes doigts}{act-ps-11}\ligneAct{12}{Le dé et les bonbons}{act-ps-12}\ligneAct{13}{Autant que sur la carte}{act-ps-13}\ligneAct{14}{Les doigts des deux mains}{act-ps-14}\ligneAct{15}{Des trains de deux couleurs}{act-ps-15}\ligneAct{16}{Quatre}{act-ps-16}\ligneAct{17}{Les chiffres 1, 2, 3}{act-ps-17}\ligneAct{18}{La chenille qui compte}{act-ps-18}}
\begin{activite}{7}{Un seul}{colorier}{Colorie le rond sous chaque boîte où il y a {\otBold\color{otNavy}un seul} objet.}{Percevoir globalement une petite quantité : un, deux.}{act-ps-7}
\matFigurePage{ps-07}
\end{activite}
\begin{activite}{8}{Deux}{colorier}{Colorie le rond sous chaque boîte où il y a {\otBold\color{otNavy}deux} objets.}{Reconnaître des collections de deux objets de tailles et de dispositions différentes.}{act-ps-8}
\matFigurePage{ps-08}
\end{activite}
\begin{activite}{9}{Trois}{colorier}{Colorie le rond sous chaque boîte où il y a {\otBold\color{otNavy}trois} objets.}{Reconnaître des collections de trois objets, quelle que soit leur disposition.}{act-ps-9}
\matFigurePage{ps-09}
\end{activite}
\begin{activite}{10}{Un de plus}{dessiner}{Dans chaque boîte de droite, dessine {\otBold\color{otNavy}un rond de plus} que dans la boîte de gauche.}{Réaliser une collection contenant un objet de plus qu’une collection donnée.}{act-ps-10}
\matFigurePage{ps-10}
\end{activite}
\begin{activite}{11}{Mes doigts}{relier}{Relie chaque main à la boîte qui a {\otBold\color{otNavy}autant} de pommes que de doigts levés.}{Associer une quantité à sa représentation avec les doigts.}{act-ps-11}
\matFigurePage{ps-11}
\end{activite}
\begin{activite}{12}{Le dé et les bonbons}{relier}{Relie chaque dé à la boîte qui a {\otBold\color{otNavy}autant} de bonbons que de points.}{Associer une constellation de points à la quantité correspondante.}{act-ps-12}
\matFigurePage{ps-12}
\end{activite}
\begin{activite}{13}{Autant que sur la carte}{dessiner}{Dessine dans chaque boîte {\otBold\color{otNavy}autant} de ronds qu’il y a de points sur la carte.}{Constituer une collection de même quantité qu’une représentation de points.}{act-ps-13}
\matFigurePage{ps-13}
\end{activite}
\begin{activite}{14}{Les doigts des deux mains}{entourer,dire}{Combien de doigts sont levés en tout ? Entoure le dé qui montre {\otBold\color{otNavy}autant} de points.}{Composer deux et trois avec les doigts des deux mains.}{act-ps-14}
\matFigurePage{ps-14}
\end{activite}
\begin{activite}{15}{Des trains de deux couleurs}{colorier}{Colorie les wagons de chaque train en rouge et en bleu. Utilise les deux couleurs dans chaque train.}{Décomposer deux et trois : deux, c’est un et un ; trois, c’est un et deux.}{act-ps-15}
\matFigurePage{ps-15}
\end{activite}
\begin{activite}{16}{Quatre}{colorier}{Colorie le rond sous chaque boîte où il y a {\otBold\color{otNavy}quatre} objets.}{Dénombrer une collection jusqu’à quatre, en déplaçant ou en pointant les objets.}{act-ps-16}
\matFigurePage{ps-16}
\end{activite}
\begin{activite}{17}{Les chiffres 1, 2, 3}{relier}{Relie chaque chiffre à la boîte qui a {\otBold\color{otNavy}autant} d’étoiles.}{Associer une écriture chiffrée à une quantité (jusqu’à trois).}{act-ps-17}
\matFigurePage{ps-17}
\end{activite}
\begin{activite}{18}{La chenille qui compte}{dire,colorier}{Touche chaque rond de la chenille en disant : un, deux, trois… jusqu’à six. Puis colorie-la.}{Réciter la comptine numérique de un à six, de façon ordonnée, en partant de un.}{act-ps-18}
\matFigurePage{ps-18}
\end{activite}
\partieMat{3}{Petits problèmes}{Ajouter, enlever : se représenter l’action.}{\pic[scale=2.6,mat/c=matBrun] at (15,8) {ours};\pic[scale=2.4,mat/c=matOrange] at (18.2,6.8) {valise};\pic[scale=1.5,mat/c=matJaune] at (18.6,9.8) {ours};}{Les problèmes sont d’abord joués avec du matériel : on ajoute, on enlève, on cache. La fiche vient après, quand la situation est comprise. L’enfant peut s’aider de ses doigts.}{\ligneAct{19}{Les ours dans la boîte}{act-ps-19}\ligneAct{20}{La boîte de crayons}{act-ps-20}\ligneAct{21}{Les oiseaux s’envolent}{act-ps-21}\ligneAct{22}{Une coccinelle arrive}{act-ps-22}}
\begin{activite}{19}{Les ours dans la boîte}{entourer}{Il y a deux ours dans la boîte. On en met encore un. Combien y a-t-il d’ours dans la boîte ? Entoure le bon dé.}{Résoudre un petit problème d’ajout en se représentant l’action.}{act-ps-19}
\matFigurePage{ps-19}
\end{activite}
\begin{activite}{20}{La boîte de crayons}{entourer}{Il y avait quatre crayons dans la boîte. On en a sorti deux. Combien en reste-t-il dans la boîte ? Entoure la bonne main.}{Résoudre un petit problème de retrait en se représentant l’action.}{act-ps-20}
\matFigurePage{ps-20}
\end{activite}
\begin{activite}{21}{Les oiseaux s’envolent}{entourer}{Il y avait trois oiseaux sur la branche. Un oiseau s’est envolé. Combien en reste-t-il sur la branche ? Entoure la bonne carte.}{Résoudre un petit problème de retrait en percevant la solution.}{act-ps-21}
\matFigurePage{ps-21}
\end{activite}
\begin{activite}{22}{Une coccinelle arrive}{entourer}{Une coccinelle est sur la feuille. Une autre arrive. Combien y aura-t-il de coccinelles sur la feuille ? Entoure la bonne main.}{Résoudre un petit problème d’ajout : un et encore un.}{act-ps-22}
\matFigurePage{ps-22}
\end{activite}
\partieMat{4}{Les formes}{Reconnaître, trier, encastrer, assembler.}{\draw[mat contour,fill=matRouge] (15,8) circle (1.2);\draw[mat contour,fill=matJaune] (17.565,8.204) -- (19.196,8.965) -- (18.435,10.596) -- (16.804,9.835) -- cycle;\draw[mat contour,fill=otGreen] (16.815,4.533) -- (19.11,5.138) -- (17.392,6.645) -- cycle;}{L’enfant reconnaît et trie des objets selon leur forme, comprend qu’une forme ne change pas quand on la tourne et reproduit de petits assemblages (au plus quatre pièces). L’adulte emploie les bons mots — disque, carré, triangle — sans les exiger.}{\ligneAct{23}{Pareil que le modèle}{act-ps-23}\ligneAct{24}{Ranger les formes}{act-ps-24}\ligneAct{25}{Chaque pièce dans son trou}{act-ps-25}\ligneAct{26}{Les deux moitiés}{act-ps-26}\ligneAct{27}{La pièce qui manque}{act-ps-27}\ligneAct{28}{Le carrelage}{act-ps-28}}
\begin{activite}{23}{Pareil que le modèle}{entourer}{Entoure toutes les formes qui ont la {\otBold\color{otNavy}même forme} que le modèle.}{Reconnaître des objets de même forme, quelles que soient leur couleur et leur taille.}{act-ps-23}
\matFigurePage{ps-23}
\end{activite}
\begin{activite}{24}{Ranger les formes}{relier}{Relie chaque forme à la boîte qui a la {\otBold\color{otNavy}même forme}.}{Trier et classer des formes planes selon leur forme, malgré la couleur, la taille et l’orientation.}{act-ps-24}
\matFigurePage{ps-24}
\end{activite}
\begin{activite}{25}{Chaque pièce dans son trou}{relier}{Relie chaque pièce au trou où elle peut entrer.}{Percevoir qu’une forme reste la même quand on la tourne ; encastrer des objets.}{act-ps-25}
\matFigurePage{ps-25}
\end{activite}
\begin{activite}{26}{Les deux moitiés}{relier}{Relie chaque moitié d’image à l’autre moitié.}{Reconstituer une image coupée en deux (assemblage de deux pièces).}{act-ps-26}
\matFigurePage{ps-26}
\end{activite}
\begin{activite}{27}{La pièce qui manque}{entourer}{Il manque une pièce au puzzle de l’ours. Entoure la pièce qui manque.}{Compléter un assemblage de quatre pièces à partir d’un modèle.}{act-ps-27}
\matFigurePage{ps-27}
\end{activite}
\begin{activite}{28}{Le carrelage}{colorier}{Colorie le grand carrelage comme le modèle.}{Reproduire un assemblage de quatre formes à partir d’un modèle (à une autre taille).}{act-ps-28}
\matFigurePage{ps-28}
\end{activite}
\partieMat{5}{Long et court}{Comparer des longueurs.}{\pic[scale=1,mat/c=matVert,mat/l=6] at (12.4,9.4) {serpent};\pic[scale=1,mat/c=matOrange,mat/l=4.6] at (13.6,7) {serpent};}{Avant quatre ans, on compare des longueurs très différentes d’un coup d’œil, puis des longueurs proches en mettant les objets à la même origine.}{\ligneAct{29}{Le serpent le plus long}{act-ps-29}\ligneAct{30}{Le crayon le plus long}{act-ps-30}\ligneAct{31}{Même longueur}{act-ps-31}\ligneAct{32}{La plus longue, la plus courte}{act-ps-32}}
\begin{activite}{29}{Le serpent le plus long}{entourer}{Entoure le serpent le plus {\otBold\color{otNavy}long}.}{Comparer visuellement deux longueurs très différentes.}{act-ps-29}
\matFigurePage{ps-29}
\end{activite}
\begin{activite}{30}{Le crayon le plus long}{colorier}{Colorie le crayon le plus {\otBold\color{otNavy}long}.}{Comparer des longueurs quand les objets partent du même point.}{act-ps-30}
\matFigurePage{ps-30}
\end{activite}
\begin{activite}{31}{Même longueur}{relier}{Relie chaque ruban de gauche au ruban de droite qui a la {\otBold\color{otNavy}même longueur}.}{Reconnaître un objet de même longueur qu’un objet donné.}{act-ps-31}
\matFigurePage{ps-31}
\end{activite}
\begin{activite}{32}{La plus longue, la plus courte}{colorier}{Colorie la brique la plus longue en rouge et la plus courte en bleu.}{Comparer des longueurs ; préparer le rangement du plus long au plus court.}{act-ps-32}
\matFigurePage{ps-32}
\end{activite}
\partieMat{6}{Les motifs}{Copier, compléter, deviner ce qui est caché.}{\pic[scale=1.4,mat/c=matRouge] at (12,7.8) {perle};\pic[scale=1.4,mat/c=matBleu] at (13.1,8.3) {perle};\pic[scale=1.4,mat/c=matRouge] at (14.2,7.8) {perle};\pic[scale=1.4,mat/c=matBleu] at (15.3,8.3) {perle};\pic[scale=1.4,mat/c=matRouge] at (16.4,7.8) {perle};\pic[scale=1.4,mat/c=matBleu] at (17.5,8.3) {perle};\pic[scale=1.4,mat/c=matRouge] at (18.6,7.8) {perle};}{L’enfant mémorise et reproduit des motifs répétitifs très simples, souvent une alternance de deux éléments, et complète un motif dont une partie est cachée.}{\ligneAct{33}{Le collier}{act-ps-33}\ligneAct{34}{La même rangée}{act-ps-34}\ligneAct{35}{Ce qui est caché}{act-ps-35}\ligneAct{36}{Le train à copier}{act-ps-36}}
\begin{activite}{33}{Le collier}{colorier}{Colorie les perles qui manquent : le collier fait rouge, bleu, rouge, bleu…}{Compléter un motif répétitif très simple (alternance de deux couleurs).}{act-ps-33}
\matFigurePage{ps-33}
\end{activite}
\begin{activite}{34}{La même rangée}{colorier}{Colorie la rangée du bas comme celle du haut.}{Reproduire à l’identique un motif répétitif.}{act-ps-34}
\matFigurePage{ps-34}
\end{activite}
\begin{activite}{35}{Ce qui est caché}{entourer}{Quelle forme est cachée sous la carte bleue ? Entoure-la en bas.}{Compléter un motif répétitif : anticiper l’élément caché.}{act-ps-35}
\matFigurePage{ps-35}
\end{activite}
\begin{activite}{36}{Le train à copier}{colorier}{Colorie le train du bas comme celui du haut.}{Reproduire à l’identique un motif répétitif.}{act-ps-36}
\matFigurePage{ps-36}
\end{activite}
\partieMat{7}{Se repérer}{Dans, dehors, sur, sous, devant, derrière.}{\pic[scale=2.2] at (14.6,7) {souris};\pic[scale=2] at (18.4,6.6) {fromage};\pic[scale=2] at (16.8,9.6) {arbre};}{Ces mots de l’espace s’apprennent d’abord avec le corps et les objets de la classe. La fiche permet ensuite de les retrouver sur une image et de suivre un parcours.}{\ligneAct{37}{Dans le bocal}{act-ps-37}\ligneAct{38}{Sur ou sous la table}{act-ps-38}\ligneAct{39}{Devant ou derrière}{act-ps-39}\ligneAct{40}{Le chemin de la souris}{act-ps-40}}
\begin{activite}{37}{Dans le bocal}{colorier}{Colorie les bonbons qui sont {\otBold\color{otNavy}dans} le bocal.}{Comprendre et utiliser les mots « dans » et « dehors ».}{act-ps-37}
\matFigurePage{ps-37}
\end{activite}
\begin{activite}{38}{Sur ou sous la table}{entourer}{Entoure le chat qui est {\otBold\color{otNavy}sous} la table.}{Comprendre et utiliser les mots « sur » et « sous ».}{act-ps-38}
\matFigurePage{ps-38}
\end{activite}
\begin{activite}{39}{Devant ou derrière}{entourer}{Entoure le lapin qui est {\otBold\color{otNavy}derrière} l’arbre.}{Comprendre et utiliser les mots « devant » et « derrière ».}{act-ps-39}
\matFigurePage{ps-39}
\end{activite}
\begin{activite}{40}{Le chemin de la souris}{tracer}{Suis le chemin avec ton doigt, puis trace-le avec un crayon, de la souris jusqu’au fromage.}{Suivre et reproduire un parcours simple.}{act-ps-40}
\matFigurePage{ps-40}
\end{activite}
